\begin{statementbox}
\begin{center}\begin{tikzpicture}[font=\small,>=stealth,line join=round]
\draw[gray!40,rounded corners=3pt] (0,0) rectangle (14.2,7.8);\node at (7.1,7.2) {关于原点成对出现；原点本身不移动};
\begin{scope}[shift={(3.8,3.5)},x=1.25cm,y=1.1cm]\draw[->,gray] (-1.8,0)--(1.8,0) node[right]{$x$};\draw[->,gray] (0,-2.0)--(0,2.6) node[above]{$y$};\draw[blue,thick,domain=-1.25:1.25,samples=150] plot(\x,{(\x)^3});\draw[gray,dashed] (-1,-1)--(1,1);\fill[orange!85!black] (1,1) circle(2.2pt) node[above right=5pt] {$(1,1)$};\fill[orange!85!black] (-1,-1) circle(2.2pt) node[below left=5pt] {$(-1,-1)$};\fill[orange!85!black] (0,0) circle(2.2pt) node[below right=5pt] {$O(0,0)$};\end{scope}\node[align=left,anchor=west] at (7.8,4.4){$f(x)=x^3$\\[11pt]$(x,y)\longleftrightarrow(-x,-y)$\\[11pt]$x=0$ 时，横坐标不变\\[11pt]$f(0)=-f(0)$，只能为 $0$};
\node[text width=12.8cm,align=center] at (7.1,.65) {图取 $x^3$；一般结论还必须先确认 $0$ 在定义域内。};
\end{tikzpicture}\end{center}


  \textbf{\textcolor{CStatement}{这条结论解决什么问题？}}

  已知函数是奇函数，却不知道完整的函数表达式时，能否直接确定它在 $x=0$ 处的值？
  可以，但前提是这个函数在 $x=0$ 处有定义。遇到含未知数的表达式，还能利用这一点快速列方程。

  \medskip
  \textbf{\textcolor{CStatement}{成立条件}}
  \begin{description}[style=nextline, leftmargin=2.8em, labelsep=0.5em, itemsep=0.35em]
    \item[\textbf{条件一：}] $f(x)$ 是奇函数，即在定义域内成对出现的 $x$ 与 $-x$ 满足 $f(-x)=-f(x)$。
    \item[\textbf{条件二：}] $0$ 在 $f(x)$ 的定义域内，可以计算 $f(0)$。
  \end{description}

  \medskip
  \textbf{\textcolor{CStatement}{核心结论}}
  \[
    \boxed{f(0)=0}
  \]
  换成图像语言：\textbf{如果奇函数的图像碰到 $y$ 轴，那么只能在原点碰到。}

  \medskip
  \end{statementbox}
\clearpage
\begin{statementbox}
\textbf{\textcolor{CStatement}{可直接用于判断}}
  \[
    \boxed{0\text{ 有定义，且 }f(0)\ne0
    \ \Longrightarrow\ f(x)\text{ 不是奇函数}.}
  \]

  \medskip
  \textbf{\textcolor{CStatement}{不能反过来说}}
\begin{center}\begin{tikzpicture}[font=\small,>=stealth,line join=round]
\draw[gray!40,rounded corners=3pt] (0,0) rectangle (14.2,8.5);\node at (7.1,7.9) {易错二：同样经过原点，只有一幅是奇函数图像};
\node at (3.5,6.85){$y=x^3$：原点对称};\node at (10.6,6.85){$y=x^2$：仅过原点还不够};\begin{scope}[shift={(3.5,3.3)},x=1cm,y=1cm]\draw[->,gray] (-2.1,0)--(2.1,0) node[right]{$x$};\draw[->,gray] (0,-1.9)--(0,2.8) node[above]{$y$};\draw[blue,thick,domain=-1.2:1.35,samples=150] plot(\x,{(\x)^3});\fill[orange!85!black] (1,1) circle(2.2pt) node[above right=5pt] {$(1,1)$};\fill[orange!85!black] (-1,-1) circle(2.2pt) node[below left=5pt] {$(-1,-1)$};\fill[orange!85!black] (0,0) circle(2.2pt);\end{scope}\begin{scope}[shift={(10.5,3.3)},x=1cm,y=1cm]\draw[->,gray] (-2.1,0)--(2.1,0) node[right]{$x$};\draw[->,gray] (0,-1.9)--(0,2.8) node[above]{$y$};\draw[blue,thick,domain=-1.6:1.6,samples=150] plot(\x,{(\x)^2});\fill[orange!85!black] (1,1) circle(2.2pt) node[above right=5pt] {$(1,1)$};\fill[orange!85!black] (-1,1) circle(2.2pt) node[above left=5pt] {$(-1,1)$};\draw[red,fill=white] (-1,-1) circle(2.5pt);\node[red,right=5pt] at (.25,-1.2){缺少 $(-1,-1)$};\fill[orange!85!black] (0,0) circle(2.2pt);\end{scope}
\node[text width=12.8cm,align=center] at (7.1,.65) {空心红点只表示应有但缺少的对称点，不属于 $y=x^2$ 的图像。};
\end{tikzpicture}\end{center}



  仅知道 $f(0)=0$，\textbf{不能}认定函数是奇函数；例如 $f(x)=x^2$ 也满足 $f(0)=0$，却不是奇函数。

  \medskip
  \textbf{\textcolor{CStatement}{一眼检查}}

  \textbf{先看 $x=0$ 能不能代入，再用奇函数的性质；若 $0$ 不在定义域内，就不能写 $f(0)=0$。}
\end{statementbox}
